\documentclass[stu]{apa7}

\usepackage[spanish]{babel}
\usepackage[utf8]{inputenc}
\usepackage{hyperref}
\usepackage{enumitem}
\usepackage{geometry}
\usepackage{times}
\usepackage{array}
\usepackage{longtable}
\usepackage{booktabs}
\usepackage{colortbl}
\usepackage{xcolor}
\usepackage{tabularx}
\usepackage{multirow}

\geometry{margin=1in}

\title{Actividad 4 --- Fase Inicial Scrum \\ Sistema de información web contable sobre recibos de retención y movimientos de caja de la empresa minera MINCH SRL.}
\author{...}
\affiliation{...}
\course{...}
\professor{...}
\duedate{...}

\setlength{\parindent}{0.5in}

\newcolumntype{L}[1]{>{\raggedright\arraybackslash}p{#1}}
\newcolumntype{C}[1]{>{\centering\arraybackslash}p{#1}}

\begin{document}

\maketitle

\section{Contexto del proyecto}

\begin{longtable}{|L{4cm}|L{10cm}|}
\hline
\rowcolor[gray]{0.9}
\textbf{Campo} & \textbf{Valor} \\
\hline
Nombre del proyecto & SIC-MINCH --- Sistema de Información Contable para MINCH SRL. \\
\hline
Procesos a automatizar & 
\begin{itemize}[nosep]
    \item Emisión de recibos de retención con cálculo automático de descuentos (RC-IVA, IUE, IT).
    \item Registro y control de movimientos de caja (ingresos/egresos en efectivo).
    \item Gestión del libro de bancos con numeración automática por año fiscal.
    \item Consulta de estados de cuenta por cliente y proveedor.
    \item Generación de reportes financieros en PDF y Excel.
\end{itemize} \\
\hline
Alcance inicial identificado & 
\begin{itemize}[nosep]
    \item Módulo de retenciones con cálculo automático de descuentos.
    \item Módulo de libro de bancos con numeración por año fiscal (octubre -- septiembre).
    \item Módulo de caja con numeración por mes calendario.
    \item Módulo de estados de cuenta (clientes y proveedores).
    \item Módulo de reportes (PDF y Excel).
    \item Módulo de administración (usuarios, roles, permisos, maestros).
\end{itemize} \\
\hline
Usuarios clave & Administrador del Sistema, Contador General, Auxiliar Contable, Gerente. \\
\hline
\end{longtable}

\newpage

\section{Visión del producto}

\begin{center}
\fbox{\begin{minipage}{0.9\textwidth}
\vspace{0.3cm}
\textbf{VISIÓN DEL PRODUCTO}

\vspace{0.3cm}
\textbf{Para} el área contable de la empresa minera MINCH SRL. y su personal administrativo,

\textbf{que} actualmente depende de hojas Excel y procesos manuales para gestionar retenciones, movimientos de caja y libro de bancos,

\textbf{el} SIC-MINCH (Sistema de Información Contable para MINCH SRL.)

\textbf{es} un sistema web moderno basado en Laravel y Livewire

\textbf{que} automatiza el cálculo de retenciones, la numeración correlativa de comprobantes, el registro de movimientos bancarios y de caja, y la generación de reportes financieros,

\textbf{permitiendo} reducir en un 80\% el tiempo dedicado a tareas administrativas manuales, eliminando errores de cálculo y centralizando toda la información financiera en una sola plataforma accesible desde cualquier dispositivo con conexión a internet.

\vspace{0.3cm}
\textbf{Diferenciadores clave:}
\begin{itemize}
    \item Cálculo automático de descuentos en retenciones según normativa boliviana (RC-IVA, IUE, IT).
    \item Numeración correlativa automática con bloqueo concurrente para evitar duplicados.
    \item Dashboard ejecutivo con KPIs y gráficos de tendencias en tiempo real.
    \item Exportación de reportes en PDF y Excel con un solo clic.
\end{itemize}
\vspace{0.3cm}
\end{minipage}}
\end{center}

\newpage

\section{Personas (User Personas)}

\subsection{Persona 1: Carlos --- Auxiliar Contable}

\begin{longtable}{|L{3.5cm}|L{11cm}|}
\hline
\rowcolor[gray]{0.9}
\textbf{Atributo} & \textbf{Descripción} \\
\hline
Rol & Auxiliar Contable \\
\hline
Edad & 28 años \\
\hline
Experiencia & 3 años en el área contable de MINCH SRL. \\
\hline
Necesidad principal & Registrar movimientos de caja y bancos de forma rápida y sin errores, emitir recibos de retención sin perder tiempo en cálculos manuales. \\
\hline
Dolor actual & 
\begin{itemize}[nosep]
    \item Dedica aproximadamente el 40\% de su jornada a tareas manuales en Excel.
    \item
    \item Debe calcular retenciones con calculadora, propenso a errores.
    \item Las versiones de Excel se duplican al compartirlas por correo.
    \item Pierde tiempo buscando registros históricos en archivos desorganizados.
\end{itemize} \\
\hline
Frustración principal & «No puedo confiar al 100\% en los saldos de Excel porque a veces alguien más modificó una fórmula sin avisar.» \\
\hline
\end{longtable}

\subsection{Persona 2: Carmen --- Contador General}

\begin{longtable}{|L{3.5cm}|L{11cm}|}
\hline
\rowcolor[gray]{0.9}
\textbf{Atributo} & \textbf{Descripción} \\
\hline
Rol & Contador General \\
\hline
Edad & 35 años \\
\hline
Experiencia & 8 años en contabilidad, 2 años en MINCH SRL. \\
\hline
Necesidad principal & Supervisar retenciones, conciliar libro de bancos, generar estados de cuenta para proveedores y clientes, y producir reportes financieros confiables. \\
\hline
Dolor actual & 
\begin{itemize}[nosep]
    \item No tiene visibilidad en tiempo real del estado financiero de las cuentas.
    \item La conciliación bancaria requiere revisar manualmente Excel contra extractos.
    \item Generar un reporte mensual toma horas de trabajo manual.
    \item Los estados de cuenta se elaboran uno por uno a pedido.
\end{itemize} \\
\hline
Frustración principal & «Cada vez que el gerente pide un reporte urgente, tengo que dejar todo lo que estoy haciendo para consolidar datos de varios Excel.» \\
\hline
\end{longtable}

\subsection{Persona 3: Pedro --- Gerente General}

\begin{longtable}{|L{3.5cm}|L{11cm}|}
\hline
\rowcolor[gray]{0.9}
\textbf{Atributo} & \textbf{Descripción} \\
\hline
Rol & Gerente General / Administrador \\
\hline
Edad & 45 años \\
\hline
Experiencia & 15 años en el sector minero, fundador de MINCH SRL. \\
\hline
Necesidad principal & Acceder a información financiera actualizada para tomar decisiones estratégicas sin depender de terceros. \\
\hline
Dolor actual & 
\begin{itemize}[nosep]
    \item No puede consultar saldos bancarios ni estado de retenciones sin pedir un reporte al contador.
    \item Los reportes llegan con retraso y a veces con inconsistencias.
    \item No tiene un panorama general de la salud financiera de la empresa.
\end{itemize} \\
\hline
Frustración principal & «Necesito saber en tiempo real cómo está la empresa, no esperar al cierre del mes para darme cuenta de que hubo problemas.» \\
\hline
\end{longtable}

\newpage

\section{Definición de \enquote{Listo} (DoR) y \enquote{Terminado} (DoD)}

\subsection{Definición de Ready (DoR)}

Una historia de usuario está lista para ser incluida en un Sprint cuando cumple con los siguientes criterios:

\begin{enumerate}
    \item Los criterios de aceptación son claros, verificables y están redactados.
    \item Las dependencias con otras historias están identificadas y resueltas.
    \item La historia ha sido estimada en Story Points por el equipo.
    \item La prioridad ha sido asignada por el Product Owner.
    \item El diseño técnico exploratorio ha sido realizado (spike si aplica).
    \item Los datos de prueba necesarios están disponibles o se pueden generar.
\end{enumerate}

\subsection{Definición of Done (DoD)}

Una historia de usuario se considera terminada cuando cumple con los siguientes criterios:

\begin{enumerate}
    \item El código está implementado y sigue los estándares de codificación del proyecto (Pint Laravel).
    \item Las pruebas unitarias han sido escritas y pasan correctamente.
    \item Las pruebas de integración han sido ejecutadas y pasan correctamente.
    \item El código ha sido revisado por al menos otro miembro del equipo.
    \item Las funcionalidades han sido probadas en un entorno de testing.
    \item La base de datos ha sido migrada correctamente (si aplica).
    \item La documentación técnica mínima ha sido actualizada.
    \item El Product Owner ha aprobado la funcionalidad en la Sprint Review.
\end{enumerate}

\newpage

\section{Product Backlog Inicial}

\subsection{Épicas (Epics)}

\begin{longtable}{|C{2cm}|L{5cm}|L{7cm}|}
\hline
\rowcolor[gray]{0.9}
\textbf{ID} & \textbf{Épica} & \textbf{Descripción} \\
\hline
\endhead
EP-01 & Retenciones & Gestión completa de retenciones: registro, cálculo automático de descuentos, numeración correlativa mensual, generación de recibo PDF y exportación Excel. \\
\hline
EP-02 & Libro de bancos & Gestión de movimientos bancarios por cuenta, numeración automática por año fiscal, cálculo de saldos, generación de recibos PDF. \\
\hline
EP-03 & Caja & Gestión de movimientos de caja (ingresos/egresos), numeración automática por mes calendario, cálculo de saldos, generación de recibos PDF. \\
\hline
EP-04 & Estados de cuenta & Consulta de movimientos y saldos por cliente y proveedor, generación de PDF. \\
\hline
EP-05 & Reportes y dashboard & Reportes de retenciones, caja y banco con filtros; exportación PDF y Excel; dashboard ejecutivo con KPIs y gráficos. \\
\hline
EP-06 & Administración & Gestión de usuarios, roles, permisos, proveedores, clientes, cuentas bancarias e impuestos. \\
\hline
\end{longtable}

\newpage

\subsection{Historias de usuario}

\subsubsection{Épica EP-01: Retenciones}

\begin{longtable}{|L{0.8cm}|L{4.5cm}|L{3.5cm}|L{2cm}|L{1.5cm}|}
\hline
\rowcolor[gray]{0.9}
\textbf{ID} & \textbf{Historia} & \textbf{Criterios de aceptación} & \textbf{Prioridad (MoSCoW)} & \textbf{SP} \\
\hline
\endhead
HU-01 & Como Contador, quiero registrar una retención seleccionando el tipo (Servicios/Bienes) y el monto base para que el sistema calcule automáticamente los descuentos aplicables. & 
\begin{itemize}[nosep]
    \item Formulario con tipo (S/G), proveedor, monto, fecha, descripción.
    \item Al guardar, el sistema ejecuta el cálculo automático de descuentos.
    \item Muestra mensaje de confirmación con el código generado.
    \item Valida que el monto esté entre 10.00 y 999,999,999.99.
\end{itemize} & Indispensable & 8 \\
\hline
HU-02 & Como Contador, quiero que el sistema calcule automáticamente los descuentos de una retención según las tasas configuradas para que no tenga que hacerlo manualmente. & 
\begin{itemize}[nosep]
    \item Para Servicios: aplica RC-IVA 13\% + IT 3\%.
    \item Para Bienes: aplica IUE 5\% + IT 3\%.
    \item Fórmula: total = monto / (1 - suma\_tasas/100).
    \item Cada descuento = total * tasa / 100.
\end{itemize} & Indispensable & 5 \\
\hline
HU-03 & Como Contador, quiero que el sistema asigne automáticamente un código correlativo mensual a cada retención para evitar duplicados. & 
\begin{itemize}[nosep]
    \begin{minipage}[t]{4.5cm}
    \item Código secuencial por mes calendario.
    \item Usa lockForUpdate() para evitar concurrencia.
    \item Formato: ABR-001, ABR-002, etc.
    \end{minipage} & Indispensable & 3 \\
\hline
\endhead
HU-04 & Como Contador, quiero generar un recibo de retención en PDF para entregarlo al proveedor. & 
\begin{itemize}[nosep]
    \item Botón PDF en cada fila de la lista de retenciones.
    \item El PDF incluye: datos del proveedor, montos, descuentos, código, fecha.
    \item Formato A4, márgenes 15mm.
    \item Se visualiza en el navegador para imprimir o descargar.
\end{itemize} & Indispensable & 5 \\
\hline
HU-05 & Como Contador, quiero exportar las retenciones de un mes a Excel para enviarlo a Impuestos Nacionales. & 
\begin{itemize}[nosep]
    \item Seleccionar mes y tipo (S/G).
    \item Exporta a Excel con formato profesional.
    \item Incluye: proveedor, monto, descuentos, código, fecha.
\end{itemize} & Importante & 5 \\
\hline
\end{longtable}

\newpage

\subsubsection{Épica EP-02: Libro de bancos}

\begin{longtable}{|L{0.8cm}|L{4.5cm}|L{3.5cm}|L{2cm}|L{1.5cm}|}
\hline
\rowcolor[gray]{0.9}
\textbf{ID} & \textbf{Historia} & \textbf{Criterios de aceptación} & \textbf{Prioridad (MoSCoW)} & \textbf{SP} \\
\hline
\endhead
HU-06 & Como Auxiliar Contable, quiero registrar un movimiento bancario (débito/crédito) asociado a una cuenta para llevar el control de ingresos y egresos. & 
\begin{itemize}[nosep]
    \item Seleccionar cuenta bancaria, tipo (D/C), proveedor, monto, fecha.
    \item Tipo de pago: Transferencia (T) o Cheque (CH).
    \item Si es CH, el número de cheque es obligatorio.
    \item El sistema muestra el saldo actualizado después del registro.
\end{itemize} & Indispensable & 8 \\
\hline
HU-07 & Como sistema, quiero numerar automáticamente cada transacción bancaria por año fiscal (octubre -- septiembre) para mantener la correlatividad. & 
\begin{itemize}[nosep]
    \item Número secuencial por tipo (D/C) dentro del año fiscal.
    \item Año fiscal: octubre a septiembre.
    \item Usa lockForUpdate() para evitar duplicados.
    \item Formato: D-000001, C-000001.
\end{itemize} & Indispensable & 3 \\
\hline
HU-08 & Como Auxiliar Contable, quiero ver el libro de bancos de un mes específico con el saldo corriente en cada fila para conocer el estado de la cuenta. & 
\begin{itemize}[nosep]
    \item Vista mensual por cuenta bancaria.
    \item Columnas: fecha, descripción, débito, crédito, saldo.
    \item Saldo calculado con función de ventana SUM OVER.
    \item Filtro por mes y proveedor.
\end{itemize} & Indispensable & 5 \\
\hline
HU-09 & Como Auxiliar Contable, quiero generar un recibo de banco en PDF para entregar al proveedor. & 
\begin{itemize}[nosep]
    \item Botón PDF en cada movimiento.
    \item Formato carta (215.9 x 279.4 mm), márgenes 0.
    \item Incluye monto en literal (español) usando NumberHelper.
\end{itemize} & Importante & 3 \\
\hline
\end{longtable}

\subsubsection{Épica EP-03: Caja}

\begin{longtable}{|L{0.8cm}|L{4.5cm}|L{3.5cm}|L{2cm}|L{1.5cm}|}
\hline
\rowcolor[gray]{0.9}
\textbf{ID} & \textbf{Historia} & \textbf{Criterios de aceptación} & \textbf{Prioridad (MoSCoW)} & \textbf{SP} \\
\hline
\endhead
HU-10 & Como Auxiliar Contable, quiero registrar un movimiento de caja (ingreso/egreso) para controlar el efectivo. & 
\begin{itemize}[nosep]
    \item Seleccionar tipo (D/C), proveedor, monto, fecha, descripción.
    \item El sistema asigna número de caja automático por mes.
    \item Al ser el primer movimiento del mes, calcula saldo anterior.
\end{itemize} & Indispensable & 8 \\
\hline
HU-11 & Como sistema, quiero numerar automáticamente cada movimiento de caja por mes calendario para mantener la correlatividad. & 
\begin{itemize}[nosep]
    \item Número secuencial por tipo (D/C) dentro del mes calendario.
    \item Usa lockForUpdate() para evitar duplicados.
    \item Formato: DOC-000001.
\end{itemize} & Indispensable & 3 \\
\hline
HU-12 & Como Auxiliar Contable, quiero generar un recibo de caja en PDF para entregar al proveedor. & 
\begin{itemize}[nosep]
    \item Botón PDF en cada movimiento de caja.
    \item Formato carta, márgenes 0.
    \item Incluye monto en literal y datos del proveedor.
\end{itemize} & Importante & 3 \\
\hline
\end{longtable}

\newpage

\subsubsection{Épica EP-04: Estados de cuenta}

\begin{longtable}{|L{0.8cm}|L{4.5cm}|L{3.5cm}|L{2cm}|L{1.5cm}|}
\hline
\rowcolor[gray]{0.9}
\textbf{ID} & \textbf{Historia} & \textbf{Criterios de aceptación} & \textbf{Prioridad (MoSCoW)} & \textbf{SP} \\
\hline
\endhead
HU-13 & Como Contador, quiero consultar el estado de cuenta de un proveedor con el detalle de todos sus movimientos y saldo corriente para conocer su situación financiera. & 
\begin{itemize}[nosep]
    \item Pestañas para seleccionar entre proveedores y clientes.
    \item Lista de titulares con saldo actual (débito -- crédito).
    \item Al seleccionar uno, muestra movimientos paginados con saldo corriente.
    \item Totales: suma débitos, suma créditos, saldo final.
\end{itemize} & Indispensable & 8 \\
\hline
HU-14 & Como Contador, quiero generar un PDF del estado de cuenta de un proveedor o cliente para entregárselo. & 
\begin{itemize}[nosep]
    \item Botón PDF en la vista de estado de cuenta.
    \item Formato A4, márgenes 8mm.
    \item Incluye: datos del titular, todos los movimientos, saldo corriente, totales.
\end{itemize} & Importante & 3 \\
\hline
\end{longtable}

\subsubsection{Épica EP-05: Reportes y dashboard}

\begin{longtable}{|L{0.8cm}|L{4.5cm}|L{3.5cm}|L{2cm}|L{1.5cm}|}
\hline
\rowcolor[gray]{0.9}
\textbf{ID} & \textbf{Historia} & \textbf{Criterios de aceptación} & \textbf{Prioridad (MoSCoW)} & \textbf{SP} \\
\hline
\endhead
HU-15 & Como Contador, quiero generar reportes de retenciones, caja o banco con filtros por fecha y tipo para analizar la información financiera. & 
\begin{itemize}[nosep]
    \item Seleccionar tipo de reporte (retenciones, caja, banco).
    \item Filtros: rango de fechas, tipo de movimiento, cuenta (si aplica).
    \item Tabla con resultados y saldo corriente.
    \item Exportable a PDF y Excel.
\end{itemize} & Indispensable & 8 \\
\hline
HU-16 & Como Gerente, quiero ver un dashboard ejecutivo con KPIs y gráficos para conocer el estado financiero de la empresa en tiempo real. & 
\begin{itemize}[nosep]
    \item Tarjetas KPI: total proveedores, movimientos, retenciones del mes.
    \item Saldo de caja y saldo bancario consolidado.
    \item Gráficos: tendencia 12 meses, balance por cuenta, retenciones S vs G.
    \item Tablas de movimientos y retenciones recientes.
\end{itemize} & Importante & 8 \\
\hline
\end{longtable}

\subsubsection{Épica EP-06: Administración}

\begin{longtable}{|L{0.8cm}|L{4.5cm}|L{3.5cm}|L{2cm}|L{1.5cm}|}
\hline
\rowcolor[gray]{0.9}
\textbf{ID} & \textbf{Historia} & \textbf{Criterios de aceptación} & \textbf{Prioridad (MoSCoW)} & \textbf{SP} \\
\hline
\endhead
HU-17 & Como Administrador, quiero gestionar usuarios del sistema (CRUD) para controlar quién accede al sistema. & 
\begin{itemize}[nosep]
    \item CRUD completo: crear, editar, eliminar, listar.
    \item Campos: nombre, apellido, email, CI, teléfono.
    \item Asignación de roles.
    \item Contraseña por defecto: número de CI.
\end{itemize} & Indispensable & 5 \\
\hline
HU-18 & Como Administrador, quiero gestionar roles y asignar permisos a cada rol para controlar el acceso a las funcionalidades. & 
\begin{itemize}[nosep]
    \item CRUD de roles.
    \item Lista de 35 permisos disponibles.
    \item Asignar/revocar permisos con sincronización.
    \item El rol Administrador está protegido.
\end{itemize} & Indispensable & 5 \\
\hline
HU-19 & Como Contador, quiero gestionar proveedores (CRUD) para tenerlos registrados y poder asociarlos a retenciones y movimientos. & 
\begin{itemize}[nosep]
    \item CRUD completo.
    \item Búsqueda por CI o nombre.
    \item Los datos se almacenan en Person + Supplier.
    \item El CI debe ser único.
\end{itemize} & Indispensable & 3 \\
\hline
HU-20 & Como Contador, quiero gestionar clientes (CRUD) para tenerlos registrados y poder consultar sus estados de cuenta. & 
\begin{itemize}[nosep]
    \item CRUD completo.
    \item Posibilidad de adjuntar archivo (contrato).
    \item Asociación opcional a cooperativa.
\end{itemize} & Indispensable & 3 \\
\hline
HU-21 & Como Contador, quiero gestionar cuentas bancarias (CRUD) para que los movimientos se asocien a la cuenta correcta. & 
\begin{itemize}[nosep]
    \item CRUD completo.
    \item Campos: nombre, número de cuenta, iniciales, moneda (BOB/USD/EUR).
\end{itemize} & Indispensable & 3 \\
\hline
HU-22 & Como Contador, quiero gestionar impuestos/tasas (CRUD) para configurar los porcentajes de retención. & 
\begin{itemize}[nosep]
    \item CRUD completo.
    \item Campos: nombre, iniciales, porcentaje, tipo (S/G/A).
    \item Seed inicial: RC-IVA 13\%, IUE 5\%, IT 3\%.
\end{itemize} & Indispensable & 3 \\
\hline
\end{longtable}

\newpage

\subsection{Resumen del Product Backlog}

\begin{longtable}{|L{2cm}|L{3cm}|L{3cm}|L{3cm}|C{2cm}|}
\hline
\rowcolor[gray]{0.9}
\textbf{Épica} & \textbf{Indispensable} & \textbf{Importante} & \textbf{Deseable} & \textbf{Total SP} \\
\hline
\endhead
EP-01: Retenciones & HU-01 (8), HU-02 (5), HU-03 (3), HU-04 (5) & HU-05 (5) & --- & 26 \\
\hline
EP-02: Libro de bancos & HU-06 (8), HU-07 (3), HU-08 (5) & HU-09 (3) & --- & 19 \\
\hline
EP-03: Caja & HU-10 (8), HU-11 (3) & HU-12 (3) & --- & 14 \\
\hline
EP-04: Estados de cuenta & HU-13 (8) & HU-14 (3) & --- & 11 \\
\hline
EP-05: Reportes y dashboard & HU-15 (8) & HU-16 (8) & --- & 16 \\
\hline
EP-06: Administración & HU-17 (5), HU-18 (5), HU-19 (3), HU-20 (3), HU-21 (3), HU-22 (3) & --- & --- & 22 \\
\hline
\rowcolor[gray]{0.9}
\textbf{Totales} & \textbf{17 historias} & \textbf{5 historias} & \textbf{0} & \textbf{108 SP} \\
\hline
\end{longtable}

\newpage

\section{Sprint Planning --- Sprint 1}

\subsection{Datos del Sprint}

\begin{longtable}{|L{4cm}|L{10cm}|}
\hline
\rowcolor[gray]{0.9}
\textbf{Campo} & \textbf{Valor} \\
\hline
Duración del Sprint & 2 semanas (10 días hábiles) \\
\hline
Equipo & 4 desarrolladores + 1 Product Owner + 1 Scrum Master \\
\hline
Capacidad del equipo & 60 SP \newline (4 personas × 10 días × 0.75 factor enfoque = 30 días-hombre ajustados) \\
\hline
Objetivo del Sprint 1 & \textbf{«Tener la funcionalidad básica de administración del sistema, gestión de maestros y retenciones con cálculo automático, permitiendo al Contador registrar una retención completa y generar su recibo PDF.»} \\
\hline
\end{longtable}

\subsection{Historias seleccionadas}

\begin{longtable}{|C{1cm}|L{4cm}|L{2cm}|L{3cm}|L{2cm}|}
\hline
\rowcolor[gray]{0.9}
\textbf{ID} & \textbf{Historia} & \textbf{Épica} & \textbf{Descripción breve} & \textbf{SP} \\
\hline
\endhead
HU-21 & Gestionar cuentas bancarias & EP-06 & CRUD de cuentas bancarias (maestro) & 3 \\
\hline
HU-22 & Gestionar impuestos/tasas & EP-06 & CRUD de impuestos con tasas (maestro) & 3 \\
\hline
HU-19 & Gestionar proveedores & EP-06 & CRUD de proveedores (maestro) & 3 \\
\hline
HU-17 & Gestionar usuarios & EP-06 & CRUD de usuarios del sistema & 5 \\
\hline
HU-18 & Gestionar roles y permisos & EP-06 & Asignación de permisos a roles & 5 \\
\hline
HU-01 & Registrar retención con cálculo automático & EP-01 & Formulario + cálculo de descuentos & 8 \\
\hline
HU-02 & Calcular descuentos automáticamente & EP-01 & Lógica de negocio de retenciones & 5 \\
\hline
HU-03 & Numerar retención automáticamente & EP-01 & Código correlativo mensual & 3 \\
\hline
HU-04 & Generar recibo de retención PDF & EP-01 & PDF con mPDF & 5 \\
\hline
\rowcolor[gray]{0.9}
\textbf{Total} & \textbf{9 historias} & & & \textbf{40 SP} \\
\hline
\end{longtable}

\subsection{Tareas de descomposición}

\subsubsection{HU-01: Registrar retención con cálculo automático (8 SP)}

\begin{enumerate}
    \item Crear tabla retentions (migración) con campos: id, type (S/G), amount, date, description, summary, code, status, supplier\_id, user\_id.
    \item Crear tabla discounts (migración) con campos: id, amount, taxe\_id, retention\_id.
    \item Implementar modelo Retention con relaciones: belongsTo Supplier, hasMany Discount.
    \item Implementar modelo Discount con relaciones: belongsTo Retention, belongsTo Taxe.
    \item Desarrollar Livewire component RetentionComponentForm con formulario de creación/edición.
    \item Integrar búsqueda de proveedores (SupplierSearch) en el formulario.
    \item Implementar validaciones: monto entre 10.00 y 999,999,999.99, tipo obligatorio, proveedor obligatorio.
    \item Implementar método syncDiscounts() que ejecuta HU-02 y HU-03 al guardar.
    \item Desarrollar RetentionComponent con lista paginada, filtro por mes y tipo.
    \item Pruebas unitarias de creación de retención y validaciones.
\end{enumerate}

\subsubsection{HU-02: Calcular descuentos automáticamente (5 SP)}

\begin{enumerate}
    \item Implementar lógica de cálculo: total = amount / (1 - sum(tax\_rates)/100).
    \item Calcular descuento individual: discount = total * tax\_rate / 100.
    \item Crear registros Discount por cada impuesto aplicable según el tipo (S/G).
    \item Para Servicios: RC-IVA (13\%) + IT (3\%).
    \item Para Bienes: IUE (5\%) + IT (3\%).
    \item Pruebas unitarias con diferentes montos y tipos.
\end{enumerate}

\subsubsection{HU-03: Numerar retención automáticamente (3 SP)}

\begin{enumerate}
    \item Implementar método nextRetentionCode() en RetentionComponentForm o servicio.
    \item Consultar el último código del mes actual con lockForUpdate().
    \item Asignar código = último + 1, formatear como MES-XXX.
    \item Pruebas unitarias de concurrencia y secuencialidad.
\end{enumerate}

\subsubsection{HU-04: Generar recibo de retención PDF (5 SP)}

\begin{enumerate}
    \item Crear vista Blade PDF/retentionForm.blade.php con el formato del recibo.
    \item Implementar método PdfController::retentionForm().
    \item Configurar mPDF con fuente UTF-8, tamaño A4, márgenes 15mm.
    \item Mostrar datos: proveedor, monto, descuentos, total, código, fecha en español.
    \item Usar NumberHelper::toLiteral() para el monto en literal.
    \item Enlace en RetentionComponent para generar PDF de cada retención.
\end{enumerate}

\subsection{Criterios de salida del Sprint 1}

\begin{itemize}
    \item El Contador puede registrar una retención completa con cálculo automático de descuentos.
    \item El sistema asigna el código correlativo mensual automáticamente.
    \item El Contador puede generar y descargar el recibo de retención en PDF.
    \item El Administrador puede gestionar usuarios, roles, permisos, proveedores, cuentas bancarias e impuestos.
    \item Todas las historias cumplen la Definición de Terminado (DoD).
    \item El Product Owner aprueba el incremento en la Sprint Review.
\end{itemize}

\end{document}
